\section{Introduction}
\label{sec:introduction}

\subsection{Clinical setting and study premise}
\label{sec:introduction:clinical}

Pancreatic neuroendocrine tumours (\pNET{}s) are uncommon neoplasms with
heterogeneous clinical behaviour. Liver metastasis is particularly relevant
because it marks advanced disease and is associated with substantially different
management and prognosis. This report addresses two related but distinct clinical
prediction tasks in registry data. The \emph{diagnostic task} estimates whether
liver metastasis is recorded at the time of the initial cancer diagnosis. The
\emph{prognostic task} estimates overall survival among the subgroup of patients
who already have liver metastasis at diagnosis. The diagnostic outcome therefore
does not represent the later development of metastasis, and the prognostic cohort
does not represent all patients with \pNET{}.

The starting point is the multicentre study of \citet{Li2023}, which developed
diagnostic and prognostic nomograms from Surveillance, Epidemiology, and End
Results (SEER) data and evaluated them in additional hospital cohorts. SEER is a
population-based cancer surveillance programme maintained by the US National
Cancer Institute \citep{NCISEER}. Reconstructing the published workflows on the
available SEER extract establishes clinically interpretable reference analyses,
exposes differences between the published description and an executable
pipeline, and provides a common baseline for the subsequent experiments.
Reproduction in this report means reproducing the analytical logic as closely as
the available variables and methodological description allow; it does not imply
access to the original patient sample or the external validation cohorts used by
the reference study. A separate 30-patient hospital cohort from
Universit\"atsklinikum Erlangen (University Hospital Erlangen, hereafter
UK--Erlangen) was subsequently made available. It is used to evaluate, without
refitting, the diagnostic and prognostic models trained during the SEER-based
reproduction of the reference study. This cohort is independent of the reference
study's hospital cohorts and is not presented as their replication.

\subsection{Scientific motivation}
\label{sec:introduction:motivation}

Modern tabular foundation models offer more than an additional prediction
algorithm. A pretrained model may be used directly for in-context prediction,
as a representation model that maps patients into a dense embedding space, or
as a generator of synthetic tabular records. TabPFN provides a prominent example
of this paradigm for small tabular datasets \citep{Hollmann2025}. Each role must
be assessed differently. When the model is used for prediction, it must
generalise to real patients who were not used for fitting. When it is used to
create embeddings, the resulting representation should retain information that
is useful for the clinical task. When it is used to generate synthetic records,
the generated cohort should reproduce relevant clinical structure, but it must
not be treated as independent patient evidence.

Consequently, a higher score from one model is not sufficient to establish that
a transformation is faithful or that synthetic augmentation is useful. The
evaluation must also consider probability error, calibration, dataset
similarity, subgroup support and the conditions under
which uncertainty was estimated. Classical statistical and machine-learning
models remain essential comparators because they distinguish gains associated
with a foundation model from gains attributable to the evaluation protocol,
model capacity or additional training records.

The central research question is therefore:
\begin{quote}
\itshape
How reliably can diagnostic and prognostic analyses for liver-metastatic
\pNET{} be reproduced from the available SEER data, and what is gained or lost
when the real clinical variables are replaced or supplemented by TabPFN-derived
representations, synthetic records, MLP models and other tabular
foundation models?
\end{quote}

\subsection{Aim and objectives}
\label{sec:introduction:objectives}

The overall aim is to evaluate tabular foundation-model methods as predictors,
representations and data generators within a traceable clinical-registry
workflow. The work is organised around nine objectives:

\begin{enumerate}[label=\textbf{O\arabic*.},ref=O\arabic*]
\item \label{obj:data}\textbf{SEER data preparation.}
Document how the data were downloaded using SEER*Stat and how the cohort
eligibility criteria, outcome definitions and preprocessing steps were based on
the reference study and adapted where required to construct the diagnostic and
prognostic datasets.

\item \label{obj:reference}\textbf{Reference-study reproduction.}
Reconstruct the diagnostic and prognostic workflows of \citet{Li2023} using the
November 2023 SEER extract after preprocessing it according to the
reference-study methodology. Record all necessary methodological adaptations and
the limits of comparison because this extract is not the November 2019 SEER
dataset used by the reference study, which was no longer available through the
project's SEER*Stat access.

\item \label{obj:external}\textbf{External validation of the reference models.}
Apply, without refitting or recalibration, the diagnostic and prognostic models
trained during the SEER-based reference-study reproduction to the independently
supplied UK--Erlangen hospital cohort, and quantify discrimination, probability
error, calibration and survival separation.

\item \label{obj:transformations}\textbf{TabPFN-derived datasets.}
Use TabPFN to create dense representations of the observed patients and to learn
the joint distribution of the clinical variables and outcomes, from which
complete synthetic patient records can be sampled.

\item \label{obj:transformed-reproduction}\textbf{Reproduction on transformed data.}
Apply the diagnostic and prognostic analysis pipelines to the embedding and
complete synthetic datasets and compare these results with the real-data
reproduction under their respective evaluation designs.

\item \label{obj:similarity}\textbf{Similarity and synthetic-data fidelity.}
Quantify the extent to which the transformations preserve marginal
distributions, relationships among variables, local geometry, clinically
important categories, overall and subgroup survival patterns, and downstream
model performance.

\item \label{obj:foundation-models}\textbf{Multi-model benchmarking.}
Compare tabular foundation models, a survival-pretrained foundation model and
classical baselines on matched real-patient evaluation units using
task-appropriate discrimination, probability-error and calibration measures.

\item \label{obj:mlp}\textbf{Nonlinear prediction with MLPs.}
Assess whether multilayer-perceptron (MLP) models applied to the clinical
variables improve upon the performance achieved by the initial reconstructed
reference-study models. Use matched linear, tree-based and foundation-model
results as additional context, and explore MLP models applied to the saved TabPFN
embeddings separately.

\item \label{obj:augmentation}\textbf{Synthetic augmentation.}
For each evaluation split, use only the real training patients to generate a
synthetic pool. Add 1,000, 5,000, 10,000 or 15,000 of these synthetic records to
the same real training set, fit each model on the resulting combined dataset and
evaluate it only on the held-out real patients. Compare every augmented condition
with the same model trained on the real training set alone.
\end{enumerate}

\subsection{Analytical structure}
\label{sec:introduction:structure}

Table~\ref{tab:introduction:architecture} summarises how the experimental layers
answer the central question. Each layer examines a different aspect of the
project, so its findings must be interpreted according to its own purpose and
evaluation design. Transformation fidelity, predictive accuracy and augmentation
utility are therefore assessed separately.

\begin{table}[htbp]
\centering
\caption{Analytical structure of the project. ``Real test'' denotes evaluation
on observed patients who were excluded from fitting and, for augmentation,
excluded from synthetic-data generation.}
\label{tab:introduction:architecture}
\small
\begin{tabularx}{\textwidth}{@{}>{\RaggedRight\arraybackslash}p{0.19\textwidth}YY@{}}
\toprule
\textbf{Analytical layer} & \textbf{Principal comparison} & \textbf{Purpose} \\
\midrule
Data processing
  & Raw SEER extract to task-specific diagnostic and prognostic cohorts
  & Establishes the eligible populations, variables, outcomes and attrition. \\
Reference reproduction
  & Reconstructed nomogram workflows on real clinical variables
  & Provides interpretable clinical baselines and documents reproducibility limits. \\
External validation
  & SEER-trained reference models applied unchanged to independent UK--Erlangen patients
  & Tests whether the reference models retain predictive value in an independent cohort, while making the limitations of the small sample clear. \\
TabPFN transformations
  & Original clinical data compared with TabPFN embeddings and synthetic records containing clinical variables and outcomes
  & Tests whether predictive and structural information survives transformation. \\
Foundation-model benchmark
  & Tabular foundation models and classical baselines trained and tested on the same SEER data splits
  & Compares model families using the same patient groups and evaluation measures. \\
MLP models
  & MLPs compared with the initial reference models and with additional models trained and tested on the same SEER data splits
  & Tests whether additional nonlinear capacity improves the previously achieved performance. \\
Synthetic augmentation
  & Real-only training versus real plus 1K, 5K, 10K or 15K synthetic rows
  & Measures whether training-only generated records improve performance on held-out real patients. \\
\bottomrule
\end{tabularx}
\end{table}
\FloatBarrier

\subsection{Scope and principles of interpretation}
\label{sec:introduction:scope}

This study evaluates how well models predict outcomes from previously recorded
SEER data. It does not determine why an outcome occurred or whether a treatment
caused benefit or harm. Treatment information in the registry may be incomplete,
and some treatments may not have been known at the time when a prediction would
need to be made. The results should therefore be interpreted as an evaluation of
predictive performance, not as evidence for treatment decisions or causal
relationships.

Most comparative experiments use internal SEER validation. The reference study's
hospital cohorts were unavailable, but an independently supplied UK--Erlangen
cohort of 30 patients permits external validation of the conventional diagnostic
and prognostic models trained during the SEER-based reproduction. Only 11
patients enter the prognostic analysis. This small
evaluation supports a preliminary assessment of transportability, not precise
subgroup inference, deployment readiness or patient benefit. The foundation-
model, representation, MLP and augmentation experiments remain internally
evaluated and are not externally validated by this cohort.

The uncertainty estimates are interpreted according to what was resampled.
Intervals obtained by resampling predictions from fixed fitted models quantify
test-patient sampling variability conditional on those models. They do not
include the uncertainty that would arise from repeating model selection,
training or synthetic-data generation. Repeated folds and overlapping test sets
are likewise not treated as independent studies.

The project uses synthetic records in two ways. First, complete synthetic
datasets containing both clinical variables and outcomes are compared with the
original SEER data to examine which distributions and clinical relationships are
preserved. These generated records are not independent patients and therefore
cannot be used as external validation data. Second, the augmentation experiments
generate synthetic records using only the real patients assigned to the training
set in each evaluation split. The generated records are added only to that
training set, and model performance is measured on the held-out real patients.
Important limitations remain: grade~4 is absent from the synthetic data, and a
small proportion of generated records contain a tumour-size category that does
not agree with the generated tumour size in millimetres. Analyses based on the
saved embeddings are considered exploratory because those embeddings were
created from the full cohort before the later training and test splits were
defined.

\subsection{Organisation of the report}
\label{sec:introduction:organisation}

The report follows the analytical sequence in Table~\ref{tab:introduction:architecture}.
It first documents SEER data acquisition, cohort construction and preprocessing,
then presents the reference diagnostic and prognostic reproductions and their
evaluation in the independent UK--Erlangen cohort. Subsequent sections describe
the TabPFN transformations, reproduction on transformed data,
dataset-similarity and synthetic-fidelity analyses, and a comparison of tabular
foundation models trained and evaluated on the same SEER patient splits with the
same performance measures. The final experimental sections assess nonlinear MLP models and
synthetic augmentation in which generated records are added only to the training
data and performance is evaluated on held-out real patients. An integrated discussion
then compares the findings across experiments, distinguishes supported results
from exploratory or unsupported ones, and closes with study limitations and
conclusions. Readers who need further technical detail can consult the additional
metric definitions, source records, code, tables and figures stored with the
project files.
